% !PRINTED 340
\weberpage{340}
% running head: 340 | Elfter Abschnitt. | §. 91.
%<<CONT
je nachdem also $T$ positiv oder negativ ist, fällt der erste oder
der zweite Werth von $r$ reell aus, und im Wesentlichen bleibt
Alles wie vorher.

Es ergiebt sich hieraus der Satz:

\gesperrt{Jede reelle ternäre cubische Form mit nicht verschwindender
 Discriminante kann durch eine reelle
Transformation auf die reelle canonische Form gebracht
 werden.}

Man kann jetzt leicht die Coordinaten der Wendepunkte in
dem canonischen Coordinatensysteme angeben. Man erhält sie,
wenn man in der Gleichung
\[
y_1^3+y_2^3+y_3^3+6m\,y_1y_2y_3=0
\]
je einmal $y_1$, $y_2$, $y_3=0$ setzt. Da eine der nicht verschwindenden
Coordinaten $=1$ gesetzt werden kann, so findet man, wenn $\varepsilon$
eine imaginäre dritte Einheitswurzel ist, die folgenden Werthe
der Coordinaten $(y_1, y_2, y_3)$:
\weq{16}{\begin{array}{ccc}
(0, 1, -1), & (0, 1, -\varepsilon), & (0, 1, -\varepsilon^2), \\
(-1, 0, 1), & (-\varepsilon, 0, 1), & (-\varepsilon^2, 0, 1), \\
(1, -1, 0), & (1, -\varepsilon, 0), & (1, -\varepsilon^2, 0),
\end{array}}
und es folgt daraus, dass \gesperrt{immer drei der Inflexionspunkte
reell und sechs imaginär sind.}

In der Tabelle (16) liegen je drei in einer Zeile stehende
Punkte auf einer Geraden, nämlich auf den Geraden
\[
y_1=0, \quad y_2=0, \quad y_3=0.
\]

Es ist danach leicht, auch die übrigen Tripel von Punkten
zusammenzustellen, die auf einer geraden Linie liegen, und damit
diese zwölf Linien zu bestimmen. Man hat dabei nur immer
solche Punkte zusammenzustellen, deren Coordinaten eine verschwindende
 Determinante geben. Man kann diese zwölf Linien
auf vier Arten in Tripel anordnen, so dass in jedem Tripel alle
neun Inflexionspunkte vorkommen. Das erste dieser Tripel wird
von den Zeilen von (16) gebildet. Die drei anderen lauten so:
\weq{17}{\begin{array}{ccc}
(0, 1, -1), & (-1, 0, 1), & (1, -1, 0), \\
(0, 1, -\varepsilon), & (-\varepsilon, 0, 1), & (1, -\varepsilon, 0), \\
(0, 1, -\varepsilon^2), & (-\varepsilon^2, 0, 1), & (1, -\varepsilon^2, 0),
\end{array}}
%CONT>>
